\section{Run \taoce}

%%%%%%%%%%%%%%%%%%%%%%%%%%%%%%%%%%%%%%%%%%%%%%%%%%%%%%%%%%%%%%%%%%%%%%%%%%%%%%%%
\subsection{Initialization}
\begin{frame}
    \alert{SystemD} is used to manage \taoce services, and therefore is the first process started in container.

    \begin{description}
        \item[devcontainer] SystemD will not start any service, as they are not built yet.
        \item[tao-ce] SystemD will \begin{enumerate} 
            \item check if configuration already exists, if not it will generate it,
            \item start the reverse-proxy,
            \item start the application services, which will initialize themselves.
        \end{enumerate} 
    \end{description}
\end{frame}

%%%%%%%%%%%%%%%%%%%%%%%%%%%%%%%%%%%%%%%%%%%%%%%%%%%%%%%%%%%%%%%%%%%%%%%%%%%%%%%%
\subsection<article>{Reverse-Proxy \& Static content}
\begin{frame}<article>
    \alert{Caddy} is used as a reverse-proxy to expose applications as a single endpoint. It is also used to serve static content.
\end{frame}

%%%%%%%%%%%%%%%%%%%%%%%%%%%%%%%%%%%%%%%%%%%%%%%%%%%%%%%%%%%%%%%%%%%%%%%%%%%%%%%%
\subsection{Applications configuration}
\begin{frame}
    Configuration files are generated at first start of container with \alert{Jsonnet}.

\only<article>{
    Those configuration files can be:
    \begin{itemize}
        \item environment variables,
        \item application specific configuration,
        \item TLS certificates,
        \item \dots
    \end{itemize}
}

    To modify applications configuration:
    \begin{itemize}
        \item[\faBan] avoid to edit configuration files manually, as they will be overwritten by the next start of the container,
        \item[\faCheck] in \devcontainer, you can directly edit \href{https://github.com/tao-ce/tao-ce/tree/main/libexec/setup}{their templates}, and apply changes with \texttt{dev:setup} task.
        \item[\faCheck] in tao-ce container, you can patch their templates from \texttt{/usr/local/libexec/tao-ce/setup} directory, and rebuild the container.
    \end{itemize}
\end{frame}

%%%%%%%%%%%%%%%%%%%%%%%%%%%%%%%%%%%%%%%%%%%%%%%%%%%%%%%%%%%%%%%%%%%%%%%%%%%%%%%%
\subsection{Applications services}
\begin{frame}[fragile=singleslide]
    Each application is composed of a set of services to be run. They are registered as \alert{systemd units}, and named after the application, e.g.:
    \begin{itemize}
        \item \texttt{tao-ce.portal.backend.service}
        \item \texttt{tao-ce.portal.bootstrap.service}
        \item \texttt{tao-ce.portal.worker.service}
        \item \texttt{tao-ce.deliver.backend.service}
        \item \texttt{tao-ce.deliver.bootstrap.service}
        \item \texttt{tao-ce.deliver.worker.service}
        \item \dots
    \end{itemize}

    You can find services in \taocerepo \texttt{apps/<app>/meta/systemd} directory of each application.

\begin{onlyenv}<article>
You can list all services with the following command: \mint{bash}{systemctl -a | grep tao-ce}
\end{onlyenv}
\end{frame}


